\begin{textsample}{POS Dim 6 – ro – Score 13.00 – ro/ro\_em\_80.txt}  \label{ex:f6_pos_131}
3.3.4.2.
Quadro Organizador das \textbf{Trilhas} de Aprofundamento da Área de Ciências Humanas e Sociais Aplicadas O quadro organizador da área de CHSA visa, de modo geral, uma interligação de conhecimentos dos componentes curriculares da área ( Filosofia, Sociologia, História e Geografia ) e das demais áreas de conhecimento.
Em relação à área de CHSA, cabe destacar que cada componente tem sua identidade e traz contribuições para o desenvolvimento humano dos estudantes, desse modo para que haja o desenvolvimento das habilidades é necessário um trabalho interdisciplinar, integrado e um olhar sensível do docente no que diz respeito a considerar a realidade, as necessidades e a peculiaridade local do estudante.
Desse modo, o fazer pedagógico deverá estar alinhado às questões cotidianas fazendo parte dessa realidade social do estudante, de maneira que os saberes escolares façam sentido para esse estudante, essa realidade precisa ser considerada pela escola no dia a dia.
Nesta ótica, todas as ações devem partir das vivências do estudante, e ainda envolver os mesmos como sujeitos culturais e sociais da sua história.
O jovem, neste processo, deve encontrar soluções criativas para a transformação social de cada ser.
As intenções didática, pedagógica e metodológica visam contemplar situações e atividades educativas que aproximam o conhecimento relacionado ao cotidiano dos estudantes, em âmbito local, regional, nacional e global, de modo a inter-relacionar e orientar sua atuação em diferentes contextos como, por exemplo : exercício das práticas cotidianas, culturais, religiosos, sociais, do trabalho, da comunicação e do seu papel enquanto cidadão.
Ademais, revela -se crucial inserir processos e modelos relativos à investigação e à resolução de \textbf{problemas} que poderão existir ao mesmo tempo em que são propostas experiências educativas que possibilitem o exercício de sua cidadania e a tomada de decisões socialmente responsáveis.
Partindo desse prisma, destaca -se que o trabalho com as \textbf{trilhas} de CHSA pode ser um desafio a mais aos estudantes, uma vez que tende a contribuir para que possam refletir sobre o modelo de sociedade na qual participa como pano de fundo têm a ciência que passa a ser um instrumento aliado a interpretação de fenômenos sociais, culturais, econômico, político e estrutural.
Concomitante, também colabora na intervenção pessoal e coletiva, promovendo aprendizagem e assumindo o papel de protagonista com responsabilidade, com dedicação de quem não se deve prender em memorizar respostas, portanto, perguntar e apresentar suas possíveis dúvidas.
Sobre esta ótica, cabe salientar que refletir sobre as culturas das quais o estudante participa, em uma sociedade em que os projetos ( Sociais e Científicos ) contribuem para incentivar e fomentar as capacidades de criar, imaginar, formular hipóteses, buscar soluções e desenvolver a sensibilidade, a empatia, a resiliência e outros aspectos positivos do ser humano que precisam ser aprimorados para viver em um mundo que se apresenta tão complexo e contraditório.
Neste contexto, o trabalho com a proposta de \textbf{trilhas} privilegia processos de autoconhecimento e de troca de saberes entre educadores e educandos, por meio de interações que se dão de forma dinâmica, criativa e construtiva o tempo todo.
Além disso, o que é pesquisado, aprendido e descoberto precisa ser compartilhado entre todos os estudantes da turma, a comunidade e outras instituições.
Face ao exposto, em relação ao quadro 4, a ilustração abaixo apresenta a descrição do código alfanumérico para o trabalho pedagógico com as \textbf{trilhas} de CHSA da seguinte forma : identificação das habilidades inerentes às habilidades associadas aos eixos \textbf{estruturantes} na área de CHSA ( 2a coluna do quadro organizador ) e as habilidades referentes às competências gerais ( 3a coluna do quadro organizador ). \textbf{Figura} 1.
Código alfanumérico do quadro organizador da área de CHSA no Ensino Médio de Rondônia.
Habilidade Específica da área associada ao eixo \textbf{estruturante} Habilidade associada às competências gerais EXEMPLO : EMIFCHS01 EM = Ensino Médio. \textbf{IF} = Itinerário \textbf{Formativo}.
CHS = Específica da área de Ciências Humanas e Sociais Aplicadas. 01 = número da habilidade.
EXEMPLO : EMIFCHS01 EM = Ensino Médio \textbf{IF} = Itinerário \textbf{Formativo} CG = Competência Geral 01 = número da habilidade.
Fonte : Retirado e adaptado da Portaria no 1432/2018 ( Brasil, 2018b ).
É importante salientar que as Competências Gerais da BNCC são alinhadas com a Agenda 2030 da Organização das Nações Unidas – ONU ( BRASIL, 2015 ), à medida que reconhece o papel da educação na afirmação de valores e no estímulo de “ ações que contribuam para a transformação da sociedade, tornando -a mais humana, socialmente justa e, também, voltada para a preservação da natureza ” ( BRASIL, 2013a ).
Busca também atender à Lei de Diretrizes e Bases da Educação e a própria Constituição, que prevê o estabelecimento de conteúdos mínimos para a educação, visando assegurar a “ formação básica comum e respeito aos valores culturais e artísticos, nacionais e regionais ” ( BRASIL, 1988 ).
Portanto, as Competências Gerais visam a “ construção de conhecimentos, no desenvolvimento de habilidades e na formação de atitudes e valores ”.
Além das Competências, foi elencado, para uma delas, um conjunto de habilidades que orientam o desenvolvimento das aprendizagens essenciais preconizadas pela BNCC.
Nesta perspectiva, o quadro 4, intitulado como Organizador das \textbf{Trilhas} de Aprofundamento foi projetado de acordo com o avanço dos processos de formação dos estudantes e suas aprendizagens fundamentais que almeja alcançar, a partir de quatro eixos \textbf{estruturantes} especificados na Portaria no 1.432/2018.
Sobre o uso desses eixos, é necessário esclarecer que oportunizam essa conexão de saberes e práticas educativas com a realidade contemporânea dos estudantes, da mesma maneira que auxilia os estudantes em seus desenvolvimentos e habilidades expressivas para sua formação integral.
Por esse motivo, na arquitetura do quadro, a primeira coluna apresenta os Eixos \textbf{Estruturantes}, sendo eles : Investigação Científica ; Processos Criativos ; \textbf{Mediação} e Intervenção Sociocultural ; e \textbf{Empreendedorismo}, que deverão ser desenvolvidos de forma interdisciplinar na área e entre as áreas de conhecimento.
Já nas segunda e terceira colunas são apresentadas, respectivamente, às Habilidades associadas aos Eixos \textbf{Estruturantes} para o trabalho na área de CHSA e as Habilidades relacionadas às Competências Gerais para o trabalho de forma integrada com as demais áreas de conhecimentos.
Na quarta coluna, os Pressupostos Metodológicos representam uma proposta metodológica com aplicações de ferramentas e elementos que objetivos novos prognósticos didáticos no processo de ensino aprendizagem da \textbf{Trilha} de Aprofundamento.
Em relação aos Eixos \textbf{Estruturantes} agrupados na área CHSA, com o crescimento e uso de \textbf{tecnologia} cada vez mais presente na educação e no cotidiano do estudante, precisa inovar e descobrir novas maneiras que possam chamar e prender a atenção dos estudantes, através de novas metodologias, procurar potencializar a aprendizagem e atrair os estudantes.
De modo específico, referente aos pressupostos metodológicos, o eixo Investigação Científica estabelece como meta e destaque, o estudante no centro do processo ensino aprendizagem e, para isso, é de fundamental importância o uso das novas metodologias ativas na aplicação da prática e nas produções científicas.
Enquanto no eixo Processos Criativos, a execução de práticas educativas transpõe o pensar e o fazer criativo do estudante sobre as artes, a cultura, as mídias e as ciências aplicadas.
Já na abordagem do eixo \textbf{Mediação} e Intervenção Sociocultural requer o uso de recursos didáticos, pedagógicos, \textbf{digitais}, tecnológicos e materiais para a investigação de conhecimentos sobre questões que afetam a vida dos seres humanos e do planeta em nível local, regional, nacional e global.
Quanto ao \textbf{Empreendedorismo}, o caminho metodológico para o processo ensino aprendizagem, requer a mobilização e união de saberes que cooperam para a identificação de potenciais, desafios, interesses e aspirações pessoais dos estudantes.
Além de que, envolve análise do contexto externo social, inclusive a relação ao mundo do trabalho ; a construção de um projeto pessoal ou produtivo ; a execução de um projeto piloto para testagem ; o aperfeiçoamento do projeto executado ; o desenvolvimento ou melhoria do projeto de vida.
Nesse sentido, o processo de desenvolvimento das \textbf{trilhas} de aprofundamento deverá ser contínuo, respeitando a dinâmica da escola e atendendo às especificidades locais e regionais.
Partir da realidade e conhecimentos locais é importante, pois o indivíduo precisa de um motivo que o impulsiona a agir, sendo assim as \textbf{trilhas} também obedecem a esse movimento.
A escola representa um pedaço da sociedade e de várias culturas em um só lugar, é pensando nessa dinâmica que existe no corpo da escola, que todos os elementos que a compõem deverão ser pensados.
No que tange às atividades que serão trabalhadas dentro do espaço da escola, elas deverão ocorrer de forma integrada, tendo como incentivos o desenvolvimento da pesquisa, convivência dos estudantes, possibilitando o trabalho em equipe, sendo incentivado e promovendo o respeito entre os estudantes.
Somente via o trabalho em equipe é que as diferenças aparecerão e a possibilidades de trabalhar questões advindas dessa convivência poderá surgir.
Em suma, cabe evidenciar que cada escola carrega consigo sua especificidade, que é \textbf{resultante} do espaço que ocupa, do \textbf{perfil} socioeconômico da região, e o processo de tutoria e \textbf{mediação}, a se ser desenvolvido pelo docente com o envolvimento coletivo de estudantes e comunidade nas situações e práticas educativas realizadas, construídas e compartilhadas, visando obter um produto.

% matched lemmas: digital, empreendedorismo, estruturante, figurar, formativo, if, mediação, perfil, problema, resultante, tecnologia, trilha
\end{textsample}
